%% chapters/04-clt.tex
\chapter{The Central Limit Theorem}
\label{ch:clt}

\whereWeAre{%
The capstone of Part~I. Lévy's continuity theorem hands us a tool to convert any pointwise statement about characteristic functions into a statement about convergence in distribution. We use it to prove the CLT --- which sits behind essentially every statement in the rest of the book.%
}

\PLACEHOLDER{Sources: existing main.tex `Characteristic Functions and Convergence in Distribution' section; the proof is already written there.}

\section{Lévy's continuity theorem}
\PLACEHOLDER{Statement, intuition, proof sketch (full proof can go to chapter appendix).}

\section{The classical CLT, via characteristic functions}
\PLACEHOLDER{Statement, Taylor expansion of the CF, $\sqrt{n}$ scaling explained, conclusion via Lévy.}

\section{Why $\sqrt{n}$? Where does the $\frac{1}{2}$ come from?}
\PLACEHOLDER{Intuition box; the second moment is the first nonzero term in the CF expansion.}

\section{The CLT in pictures}
\PLACEHOLDER{Histogram of standardized means converging to the bell curve; QQ-plot version; small Colab simulation.}

\section{Brief warm-ups}
\PLACEHOLDER{3--5 short questions.}

\chapterRecap{%
\begin{itemize}
  \item \PLACEHOLDER{The CLT is a statement about $\sqrt{n}$-scaled fluctuations.}
  \item \PLACEHOLDER{Lévy's theorem is the bridge between CFs and distributional convergence.}
  \item \PLACEHOLDER{This Gaussian universality is one of two big universality phenomena in this book; the other --- Tracy--Widom --- waits in Part~V.}
\end{itemize}%
}

\section*{Exercises}
\PLACEHOLDER{8--15 longer exercises.}
